\chapter{General Conclusion and Perspectives}
\label{ch:pfe-conclusion}

\section{General Conclusion}

This thesis investigated whether a useful Software Bill of Materials can be
recovered when the available Android artifact is a compiled \APK rather than a
source repository or reproducible build environment. The problem is difficult
because dependency archives are merged into DEX, source-level coordinates and
versions are not preserved as a complete manifest, and \Reight may rename,
remove, inline, merge, or reorganize the evidence needed for component
attribution. Vulnerability analysis adds a further requirement: an advisory
mapping is only as trustworthy as the component identity and exact-version
evidence on which it depends.

The project addressed this problem through a static, evidence-gated
APK-to-\SBOM pipeline. The implemented MVP:

\begin{itemize}
  \item inspects APK archives, manifests, DEX files, resource and asset
  inventories, native libraries, and extraction-risk indicators;
  \item extracts packages, classes, methods, opcodes, literals, API
  references, dependencies, and structural features from all successfully
  parsed DEX files;
  \item constructs versioned Maven reference profiles by acquiring JAR/AAR
  releases and converting their bytecode to a comparable DEX representation;
  \item ranks component and version candidates through explainable
  multi-feature evidence;
  \item separates component-identity confidence from exact-version confidence
  and retains ambiguous alternatives;
  \item emits accepted evidence as \CycloneDX JSON with Package URLs;
  \item queries OSV while preserving confirmed, potential, and report-only
  vulnerability interpretations;
  \item validates emitted Maven rows against source dependency trees when
  auditable source evidence is available; and
  \item produces JSON, Markdown, and HTML reports that expose evidence,
  decisions, uncertainty, and limitations.
\end{itemize}

The contribution is therefore broader than a similarity score. It combines
binary evidence extraction, reference construction, component attribution,
standard serialization, vulnerability interpretation, source validation, and
experiment governance in one reproducible workflow. The project also records
why apparently promising changes were not promoted. Global opcode weighting,
broad API weighting, recency-based version selection, global multi-scale
fusion, broad priority-database shortlisting, and raw-opcode deep learning
were rejected, disabled, or isolated when their controls did not justify
production authority.

\subsection{Evaluation Synthesis}

The controlled evaluation confirms that the implemented mechanics work when
the expected components and reference evidence are known: 84 of 84 expected
Maven versions were recovered across 72 controlled fixtures. The clean
source-validated real slice confirms 47 of 47 emitted Maven versions. These
results demonstrate that static APK evidence can support exact
component/version claims under favorable conditions.

The broader real-APK evaluation gives a more realistic boundary. Across the
102 source-validated APKs that emitted Maven components, 927 of 1,837 emitted
versions were confirmed, corresponding to a 50.5\% \emph{emitted Maven
version confirmation rate}. This metric describes correctness among emitted
Maven rows; it is not full dependency recall. Four additional source-validated
APKs emitted no Maven components and remain separate matcher-recall and
emission-threshold controls.

Against the complete source dependency denominator, the system emitted 927
exact coordinate/version rows from 12,082 resolved source rows, or 7.7\%.
The shared reference database contained the exact source version for 1,231 of
those rows, or 10.2\%. The difference between controlled success, emitted-row
confirmation, reference coverage, and full-denominator recall is one of the
principal findings of the thesis. A matcher cannot identify a missing
reference, and a large reference catalog is not sufficient if candidate
generation introduces false emissions.

The source-known v2 lane confirmed 6,772 of 6,772 emitted rows when each APK's
candidate space was constrained to exact source-known coordinate/version
references. This result demonstrates recognition under a favorable
upper-bound condition. It does not measure unconstrained APK-only accuracy and
does not remove the remaining 5,238 source-gap rows.

The vulnerability demonstrations confirm the same authority rule. Exact and
source-consistent package/version evidence can support a confirmed advisory
mapping. A source-version mismatch must remain potential even when the
external vulnerability service correctly returns an advisory for the selected
version.

The ML/DL work supports report triage rather than component authority. The
MVP-safe structured-token Transformer provides promising report-only triage
evidence, while its production decision wiring remains disabled. The
raw-opcode Transformer failed to beat its simpler logistic-regression
baseline. The integration smoke tests show that AI sections can be visible
without changing selected components, versions, \CycloneDX content, or
vulnerability output.

\subsection{Practical Significance}

The MVP is useful when interpreted as an evidence-recovery and audit tool, not
as a replacement for a trusted build-generated \SBOM. It can recover strong
component claims, reveal version ambiguity, produce interoperable output,
guide vulnerability investigation, and identify unsupported regions such as
native-heavy or hidden-code applications.

Its practical value comes from making partial knowledge explicit. A report can
distinguish:

\begin{itemize}
  \item a source-confirmed exact Maven component;
  \item a strong component identity with an ambiguous version band;
  \item a report-only candidate requiring analyst review;
  \item a missing-reference or matcher-recall gap;
  \item a hash-backed native file observation; and
  \item a protected or dynamically loaded boundary outside static coverage.
\end{itemize}

This distinction is safer than emitting one precise-looking dependency list in
which supported and unsupported claims are indistinguishable.

\section{Answers to the Research Questions}

Table~\ref{tab:pfe-conclusion-rq-answers} summarizes the answers established by
the evaluation. The wording preserves the evidence classes and denominators
defined in Chapter~\ref{ch:pfe-evaluation}.

\begin{table}[H]
\centering
\caption{Answers to the research questions and their evidence boundaries.}
\label{tab:pfe-conclusion-rq-answers}
\begin{tabularx}{\textwidth}{L{0.08\textwidth}L{0.25\textwidth}X}
\toprule
\textbf{RQ} & \textbf{Answer} & \textbf{Established boundary}\\
\midrule
RQ1 & Exact Maven versions can be recovered under bounded conditions. &
Controlled fixtures recover 84/84 expected versions, and the clean real slice
confirms 47/47 emitted versions. Partial and full-denominator evidence shows
that this capability depends on reference availability, surviving bytecode,
and version separability.\\
RQ2 & Emitted-version quality varies substantially across APK families. &
Clean, well-covered APKs support strong exact claims. Partial, minified,
AndroidX/Kotlin-era, and large-tree cases expose more wrong-version, missing,
or ambiguous rows. Native-heavy and hidden-code cases expose unsupported
analysis boundaries and remain separate from Maven-version accuracy.\\
RQ3 & Reference coverage and matcher recall both constrain complete recovery.
& The shared database exactly covers 1,231/12,082 source rows, while the full
pipeline emits 927/12,082 exact rows. The source-known v2 lane demonstrates
that broader exact references enable recognition, but its constrained
candidate space is an upper bound rather than production accuracy.\\
RQ4 & Downstream outputs can preserve uncertainty and authority boundaries. &
\CycloneDX components carry confidence and ambiguity evidence; OSV findings
inherit component/source-validation status; native files remain hash-backed
observations; and AI assistance is report-only and cannot alter protected
outputs.\\
\bottomrule
\end{tabularx}
\end{table}

The combined answer is conditional rather than absolute. Static APK analysis
can produce a useful and auditable \SBOM, but trustworthy exact-version
recovery requires adequate references, separable compiled evidence, cautious
candidate generation, and downstream systems that preserve uncertainty.

\section{Future Perspectives}

Future work should extend the evidence base in an order that reduces false
claims before increasing output volume. Table~\ref{tab:pfe-future-priorities}
summarizes the recommended priorities and their promotion conditions.

\begin{table}[H]
\centering
\caption{Prioritized future work and the evidence required for promotion.}
\label{tab:pfe-future-priorities}
\begin{tabularx}{\textwidth}{L{0.09\textwidth}L{0.25\textwidth}X}
\toprule
\textbf{Priority} & \textbf{Direction} & \textbf{Required evidence}\\
\midrule
P1 & Safer large-catalog candidate retrieval & Improve identity recall and
reference coverage without reproducing priority-database false emissions;
pass clean, partial, controlled, and same-coordinate controls.\\
P1 & Exact-version separability and calibrated abstention & Use
version-distinctive evidence and calibrated margins; demonstrate fewer wrong
versions without hiding unresolved cases.\\
P2 & Broader source-observed reference coverage & Add bytecode-bearing
source-used releases through reproducible acquisition and isolated promotion
gates; keep metadata-only and unsupported artifacts explicit.\\
P2 & Scoped multi-scale R8 evidence & Evaluate class, call-chain, basic-block,
field-operation, and literal-context evidence per risk profile; prohibit
global promotion when any protected control regresses.\\
P3 & Native component provenance & Build a separate native reference corpus
and evaluation protocol for ABI-specific shared libraries, stripping, and
static linking.\\
P3 & Dynamic and protected-code observation & Combine static risk indicators
with sandboxed runtime collection, integrity checking, and explicit
provenance for loaded payloads.\\
P4 & External benchmarks and reproducibility & Evaluate on independently
curated APKs, additional ecosystems and build eras, and compatible public
benchmarks with published denominator definitions.\\
P4 & Calibrated report-only learning & Train on external, independently
reviewed labels and APK-only inputs; require calibration and strict output
equivalence before considering any authority expansion.\\
\bottomrule
\end{tabularx}
\end{table}

\subsection{Safer Candidate Retrieval}

The highest-value near-term improvement is a stronger APK-only shortlist for
large reference catalogs. Possible mechanisms include inverted indexes over
rare classes and methods, coordinate-family retrieval, approximate nearest
neighbor search over structural summaries, and multi-stage coarse-to-fine
verification. The critical gate is not shortlist recall alone. A successful
method must recover intended candidates while preventing the extra emissions
observed in broad package-overlap experiments.

\subsection{Calibrated Evidence and Version Abstention}

Current scores rank candidates but are not formal probabilities. Future work
can calibrate identity confidence, exact-version confidence, and abstention
thresholds using held-out source-validated cases. Calibration should be
coordinate- and family-aware because score distributions differ across
libraries, build eras, and optimization profiles.

Version-set prediction may be more honest than single-version selection for
feature-equivalent releases. A future calibrated set-valued method could
produce candidate sets with explicitly evaluated coverage, provided its
assumptions and validation population are stated.

\subsection{Reference and Provenance Knowledge Graph}

The current SQLite schema stores component features and declared Maven
dependencies. A provenance-aware knowledge graph could connect:

\begin{itemize}
  \item public coordinates, target variants, aliases, and relocated
  artifacts;
  \item release chronology and version families;
  \item Maven dependency relationships;
  \item source-observed usage;
  \item binary fingerprints and evidence occurrences;
  \item vulnerability affected ranges; and
  \item the decisions and controls that justified or rejected a claim.
\end{itemize}

Such a graph could improve artifact-family reasoning and explain why a
candidate was considered. It must remain secondary to observed APK evidence;
graph coherence should not become a new form of metadata-only certainty.

\subsection{Native and Dynamic Components}

Native recovery should be developed as a separate module with its own
features, references, and metrics. Useful evidence may include exported and
imported symbols, function fingerprints, strings, section structure, build
identifiers, compiler signatures, and architecture-normalized code features.

Dynamic-code work should record when and where a payload was obtained, its
hash, signature, loading mechanism, and execution context. Runtime collection
must not be merged silently with the original \APK inventory; provenance
should distinguish packaged, unpacked, and downloaded components.

\subsection{Larger and More Independent Evaluation}

Future evaluation should increase application diversity by family rather than
by convenience. Priority categories include:

\begin{itemize}
  \item modern AndroidX and Kotlin multiplatform applications;
  \item legacy Maven and JCenter-era builds;
  \item commercial or protected applications for which analysis is legally
  authorized;
  \item games and cross-platform frameworks;
  \item native-heavy and dynamically extensible applications; and
  \item multiple releases of the same application to measure temporal drift.
\end{itemize}

An independent corpus should separate dataset construction, label review, and
tool development. It should publish both emitted-row quality and full source
or binary-ground-truth denominators so that coverage and matcher precision are
not collapsed.

\subsection{AI as a Guarded Research Direction}

Learned representations remain promising for complex \Reight evidence, but
the next step should not be a larger model on the same shortcuts. A stronger
AI study requires:

\begin{itemize}
  \item more independent APK cases per class;
  \item labels reviewed independently of the features used for prediction;
  \item APK-only inputs available at deployment time;
  \item hierarchical aggregation from methods to classes, libraries, and
  applications;
  \item calibration and out-of-distribution detection;
  \item simple and strong baselines; and
  \item a proof that authoritative outputs remain unchanged until promotion
  gates pass.
\end{itemize}

The appropriate near-term role remains analyst assistance: identify unusual
evidence patterns, prioritize review, and recommend which deterministic
diagnostic should run next.
